\chapter{Nosebleeds (Epistaxis)}

\section*{Definition}
Nosebleeds (epistaxis) are bleeding from the nasal cavity. They are very common, with approximately 60 percent of people experiencing at least one episode in their lifetime. Most nosebleeds are anterior (from Kiesselbach plexus) and self-limiting.

\section*{What Happens in Your Body}
The nasal mucosa has a rich blood supply. Anterior epistaxis (90 percent of cases) arises from the Kiesselbach plexus (Little area) on the anterior nasal septum, where several arteries anastomose. Posterior epistaxis arises from the sphenopalatine artery and is more common in older adults, more severe, and harder to control.

\section*{Causes}
\begin{itemize}
\item Nasal trauma: picking the nose, blowing too hard, foreign body
\item Dry air (low humidity, winter heating, air conditioning)
\item Nasal allergies and sinusitis
\item Upper respiratory infections
\item Nasal sprays (corticosteroids, decongestants with overuse)
\item Coagulation disorders (hemophilia, von Willebrand disease, thrombocytopenia)
\item Anticoagulant or antiplatelet medications (warfarin, aspirin, clopidogrel)
\item Liver disease (impaired coagulation factor synthesis)
\item Hereditary hemorrhagic telangiectasia (Osler-Weber-Rendu syndrome)
\item Nasal tumors (benign or malignant)
\end{itemize}

\section*{When to See a Doctor}
\begin{itemize}
\item Active bleeding from one or both nostrils
\item Blood dripping down the back of the throat (posterior nosebleed)
\item Frequent or recurrent episodes
\item Bleeding that lasts more than 20 minutes
\item Heavy blood loss causing dizziness or lightheadedness
\item Bleeding associated with easy bruising or bleeding from other sites
\item History of anticoagulant use
\end{itemize}

\section*{Related Symptoms}
\begin{itemize}
\item Nasal congestion
\item Nasal trauma
\item Sinusitis
\item Allergic rhinitis
\item Coagulopathy
\item Anemia (from chronic blood loss)
\end{itemize}

\section*{Self-Care / Home Management}
\begin{itemize}
\item Sit upright and lean forward (do NOT tilt the head back)
\item Pinch the soft part of the nose firmly for 10-15 minutes
\item Breathe through the mouth
\item Apply a cold compress or ice pack to the bridge of the nose
\item Do not pick or blow the nose for 24 hours after bleeding stops
\item Use saline nasal spray to keep the nasal passages moist
\item Use a humidifier in dry environments
\end{itemize}

\section*{How Your Doctor Will Evaluate You}
\begin{itemize}
\item Anterior rhinoscopy to identify the bleeding site
\item Nasal endoscopy to evaluate posterior bleeding
\item Complete blood count (check for anemia or thrombocytopenia)
\item Coagulation profile (PT, PTT, INR)
\item Blood pressure assessment (hypertension can worsen bleeding)
\end{itemize}

\section*{Treatment Options}
\begin{itemize}
\item Direct pressure as first-line self-care
\item Nasal packing (anterior or posterior) for refractory bleeding
\item Cauterization (chemical with silver nitrate or electrical)
\item Tranexamic acid topical or oral
\item Embolization of the sphenopalatine artery for severe posterior epistaxis
\item Ligation of the sphenopalatine or anterior ethmoid artery
\item Treat underlying cause (discontinue anticoagulants, correct coagulopathy)
\end{itemize}

\section*{References}
\begin{thebibliography}{99}
\bibitem{ref1} Krulewitz NA, Fix ML. "Epistaxis." \textit{Emergency Medicine Clinics of North America}, 2019;37(1):29--39.
\bibitem{ref2} Schlosser RJ. "Epistaxis." \textit{New England Journal of Medicine}, 2009;360(8):784--789.
\bibitem{ref3} Pope L, Hobbs CGL. "Epistaxis: an update on current management." \textit{British Medical Journal}, 2005;330(7481):26--28.
\bibitem{ref4} Tunkel DE, Anne S, Payne SC, et al. "Clinical practice guideline: nosebleed (epistaxis)." \textit{Otolaryngology--Head and Neck Surgery}, 2020;162(1\_suppl):S1--S38.
\bibitem{ref5} Beck R, Sorge M, Schneider A, et al. "Current approaches to epistaxis treatment in primary and secondary care." \textit{Deutsches \"Arzteblatt International}, 2018;115(1--2):12--18.
\bibitem{ref6} Douglas R, Wormald PJ. "Update on epistaxis." \textit{Current Opinion in Otolaryngology \& Head and Neck Surgery}, 2007;15(3):180--183.

\end{thebibliography}

\clearpage
